% ============================================================
%  PRIOR ART DEFENSIVE PUBLICATION
%  PhaseMirror: A Prime-Indexed Formally Verified Governance
%  Runtime with Contractive Kernel Semantics and Human-Gated
%  Self-Modification Protocol
%
%  MultiplicityFoundation / PhaseMirror-HQ
%  Publication Date: April 25, 2026
%  Status: Defensive Prior Art -- All Rights Reserved
% ============================================================

\documentclass[11pt,letterpaper]{article}

\usepackage[margin=1in]{geometry}
\usepackage{amsmath,amssymb,amsthm}
\usepackage{mathtools}
\usepackage[hidelinks]{hyperref}
\usepackage{listings}
\usepackage{xcolor}
\usepackage{booktabs}
\usepackage{longtable}
\usepackage{array}
\usepackage{enumitem}
\usepackage{fancyhdr}
\usepackage{titlesec}
\usepackage{float}
\usepackage{microtype}
\usepackage[T1]{fontenc}
\usepackage[utf8]{inputenc}
\usepackage{lmodern}
\usepackage{mdframed}
\usepackage{tcolorbox}
\tcbuselibrary{skins,breakable}

% ── Hyperref ────────────────────────────────────────────────
\hypersetup{
  pdftitle={PhaseMirror Prior Art Defensive Publication},
  pdfauthor={MultiplicityFoundation -- Ryan Van Gelder},
}

% ── Theorem environments ────────────────────────────────────
\newtheorem{theorem}{Theorem}[section]
\newtheorem{lemma}[theorem]{Lemma}
\newtheorem{corollary}[theorem]{Corollary}
\newtheorem{definition}[theorem]{Definition}
\newtheorem{proposition}[theorem]{Proposition}
\newtheorem{gatethm}[theorem]{Gate}

\theoremstyle{remark}
\newtheorem{invariant}[theorem]{Invariant}

% ── Code listing style ──────────────────────────────────────
\definecolor{codebg}{RGB}{248,248,248}
\definecolor{codeframe}{RGB}{200,200,200}
\definecolor{kw}{RGB}{0,0,180}
\definecolor{str}{RGB}{0,128,0}
\definecolor{cmt}{RGB}{128,128,128}
\definecolor{num}{RGB}{180,0,0}

\lstdefinestyle{python}{
  language=Python,
  backgroundcolor=\color{codebg},
  frame=single,
  rulecolor=\color{codeframe},
  basicstyle=\ttfamily\footnotesize,
  keywordstyle=\color{kw}\bfseries,
  stringstyle=\color{str},
  commentstyle=\color{cmt}\itshape,
  numberstyle=\tiny\color{num},
  numbers=left,
  numbersep=6pt,
  breaklines=true,
  showstringspaces=false,
  tabsize=4,
  captionpos=b,
}

\lstdefinestyle{mlir}{
  language={},
  backgroundcolor=\color{codebg},
  frame=single,
  rulecolor=\color{codeframe},
  basicstyle=\ttfamily\footnotesize,
  keywordstyle=\color{kw}\bfseries,
  commentstyle=\color{cmt}\itshape,
  numbers=left,
  numbersep=6pt,
  breaklines=true,
  tabsize=2,
  captionpos=b,
  morekeywords={pirtm,linalg,func,module,return,f64,i64,tensor},
}

\lstdefinestyle{cpp}{
  language=C++,
  backgroundcolor=\color{codebg},
  frame=single,
  rulecolor=\color{codeframe},
  basicstyle=\ttfamily\footnotesize,
  keywordstyle=\color{kw}\bfseries,
  stringstyle=\color{str},
  commentstyle=\color{cmt}\itshape,
  numberstyle=\tiny\color{num},
  numbers=left,
  numbersep=6pt,
  breaklines=true,
  showstringspaces=false,
  tabsize=2,
  captionpos=b,
}

% ── tcolorbox styles ────────────────────────────────────────
\newtcolorbox{invariantbox}[1][]{
  colback=blue!3!white,
  colframe=blue!50!black,
  title={#1},
  fonttitle=\bfseries\small,
  breakable,
  before skip=8pt, after skip=8pt,
}
\newtcolorbox{gapbox}[1][]{
  colback=yellow!5!white,
  colframe=orange!60!black,
  title={#1},
  fonttitle=\bfseries\small,
  breakable,
  before skip=8pt, after skip=8pt,
}

% ── Page headers ────────────────────────────────────────────
\pagestyle{fancy}
\fancyhf{}
\fancyhead[L]{\small\textit{PhaseMirror Prior Art Defensive Publication}}
\fancyhead[R]{\small\textit{Citizen Gardens, April 2026}}
\fancyfoot[C]{\thepage}
\renewcommand{\headrulewidth}{0.4pt}

% ════════════════════════════════════════════════════════════
\begin{document}
% ════════════════════════════════════════════════════════════

\begin{titlepage}
\centering
\vspace*{1.5cm}
{\large \textbf{PRIOR ART DEFENSIVE PUBLICATION}}\\[1em]
{\rule{\linewidth}{1pt}}\\[0.8em]
{\LARGE \textbf{PhaseMirror}}\\[0.5em]
{\large A Prime-Indexed Formally Verified Governance Runtime\\
with Contractive Kernel Semantics and\\
Human-Gated Self-Modification Protocol}\\[0.8em]
{\rule{\linewidth}{0.5pt}}\\[1em]
{\normalsize \textit{From Sigma Kernel to Recursive Governance:\\
A Decade Roadmap with Full Technical Specification}}\\[2em]
{\normalsize
\textbf{Ryan Van Gelder}\\
Citizen Gardens (501(c)(3))\\
\texttt{github.com/MultiplicityFoundation/PhaseMirror-HQ}\\[0.4em]
\textbf{April 25, 2026}
}\\[2em]
\begin{mdframed}[backgroundcolor=gray!10,linecolor=black,linewidth=1pt]
\small\noindent
\textbf{DEFENSIVE PUBLICATION NOTICE.}
This document constitutes a defensive prior art publication.
All systems, methods, algorithms, type-theoretic constructs, protocol
specifications, and architectural patterns described herein are disclosed
for the express purpose of placing them in the public domain as prior art,
preventing the grant of patent claims covering the same subject matter to
any third party. Publication date: \textbf{April 25, 2026}.
Released under Creative Commons CC0 for prior art purposes.
\end{mdframed}
\end{titlepage}

\tableofcontents
\newpage

% ════════════════════════════════════════════════════════════
\section{Executive Summary}
% ════════════════════════════════════════════════════════════

\textbf{PhaseMirror} is an open-source AI governance runtime developed
under MultiplicityFoundation. As of April 2026, the system is in active
development at Gate~B---the completion of the Sigma Kernel facade---with
26 roadmap gates planned through 2036. This document discloses the
complete technical specification as prior art.

\subsection{Core Claims Disclosed}

The following specific technical inventions are hereby disclosed:

\begin{enumerate}[leftmargin=2em]
  \item A compile-time contractivity type system for MLIR dialects encoding
    the invariant $r(\Lambda) < 1 - \varepsilon$ as a parameterized type
    $\texttt{!pirtm.contractivity}\langle\varepsilon,\delta\rangle$, with
    a composition algebra preserving the invariant under sequential composition.

  \item A three-checkpoint gate protocol for dialect clearance with
    mathematical acceptance criteria derived from the Banach Fixed-Point
    Theorem (Day~14), the spectral small-gain theorem (Day~30), and a
    $10\times$ NumPy throughput benchmark on 512-dimensional tensors (Day~90).

  \item An 11-canonical-constraint governance taxonomy (C01--C11) with
    a four-tier totally ordered status lattice and artifact-first evidence
    contracts.

  \item A four-phase self-modification execution protocol with proposal
    binding via SHA-256 snapshot hashes, simulation clearance,
    three-gate validation (boundary/legitimacy/playbook), and a kill switch
    that is arm-capable programmatically but disarm-capable only via human
    cryptographic token.

  \item A \texttt{VERIFICATION\_LAYER\_PARAMETERS} frozenset guard preventing
    modification of the verification machinery itself, as an architectural
    response to Lobian obstruction in governance contexts.

  \item A lax functor $F : \mathcal{F}_{\text{ext}} \to \mathcal{D}_{\text{int}}$
    for importing external governance profiles (EU AI Act, NIST AI RMF)
    into internal constraint semantics via a monotone join over the status
    lattice.

  \item A corrected Jansen-Ruskai-Seiler two-term adiabatic bound with
    exponent 6.93 for PETC evolution scheduling in prime-indexed tensor
    chains.
\end{enumerate}

\subsection{System Positioning}

PhaseMirror occupies a design space that does not currently exist in the
open-source or commercial landscape: a \emph{formally verified} governance
runtime where contractivity constraints are enforced at the \emph{type
level} of the compiler IR, rather than at runtime or via external audits.
The combination of MLIR type-level enforcement with a governance constraint
taxonomy and a human-gated self-modification protocol constitutes a
novel architectural pattern.

% ════════════════════════════════════════════════════════════
\section{System Architecture}
% ════════════════════════════════════════════════════════════

\subsection{Repository Structure}

PhaseMirror-HQ is organized as a monorepo at
\texttt{github.com/MultiplicityFoundation/PhaseMirror-HQ}.

\begin{center}
\begin{tabular}{ll}
\toprule
\textbf{Path} & \textbf{Function} \\
\midrule
\texttt{packages/governance/dcgf/} & DCGF framework \\
\texttt{packages/pirtm/} & PIRTM dialect and runtime \\
\texttt{mcp/}, \texttt{mcp\_server/} & MCP server for agent integration \\
\texttt{daemon/} & Governance daemon and modification executor \\
\texttt{lean4/} & Lean 4 formal proofs \\
\texttt{rollback/} & Snapshot ledger and revert mechanism \\
\texttt{mirror-dissonance-pro/} & Phase Mirror Dissonance (PMD) evaluator \\
\texttt{src/runtime/} & \texttt{libpirtm\_runtime.cpp} \\
\bottomrule
\end{tabular}
\end{center}

\subsection{The Four-Layer Stack}

\begin{figure}[H]
\centering
\renewcommand{\arraystretch}{1.4}
\begin{tabular}{|p{11cm}|}
\hline
\textbf{Layer 4: Governance Profiles}\\
EU AI Act $\cdot$ NIST AI RMF $\cdot$ ISO 42001 $\cdot$ Custom\\
Lax functor $F: \mathcal{F}_\text{ext} \to \mathcal{D}_\text{int}$\\
\hline
\textbf{Layer 3: DCGF Framework}\\
C01--C11 constraints $\cdot$ Legitimacy schema $\cdot$ Artifact contracts $\cdot$ Playbook\\
\hline
\textbf{Layer 2: Sigma Kernel}\\
Atomic $\|$ $\pi$ $\|$ Tensor kernel composition\\
Convergence detection $\cdot$ Audit ledger $\cdot$ MCP agent surface\\
\hline
\textbf{Layer 1: PIRTM Runtime}\\
$\texttt{!pirtm.contractivity}\langle\varepsilon,\delta\rangle$ type enforcement\\
MLIR dialect $\to$ LLVM IR $\to$ \texttt{libpirtm\_runtime.so}\\
\hline
\end{tabular}
\caption{PhaseMirror four-layer architecture}
\end{figure}

% ════════════════════════════════════════════════════════════
\section{Mathematical Foundations}
% ════════════════════════════════════════════════════════════

\subsection{The PIRTM Recurrence Equation}

The core computational object is the PIRTM state recurrence:

\begin{equation}
  X_{t+1} = P\!\left(\Xi\, X_t + \Lambda\, T(X_t) + G_t\right)
  \label{eq:recurrence}
\end{equation}

\noindent where:
\begin{itemize}
  \item $X_t \in \mathbb{R}^n$ --- state vector at step $t$
  \item $\Xi \in \mathbb{R}^{n \times n}$ --- recurrence (self-coupling) matrix
  \item $\Lambda \in \mathbb{R}^{n \times n}$ --- gain/aggregation matrix
  \item $T : \mathbb{R}^n \to \mathbb{R}^n$ --- nonlinear activation;
    $T(x) = \sigma(x) = (1+e^{-x})^{-1}$, which is $\tfrac{1}{4}$-Lipschitz
  \item $G_t \in \mathbb{R}^n$ --- external input at step $t$
  \item $P$ --- projection; implemented as $\operatorname{clip}(\cdot,-1,1)$
\end{itemize}

\subsection{The Contractivity Condition}

\begin{theorem}[PIRTM Contractivity]
\label{thm:contractivity}
The map defined by \eqref{eq:recurrence} is a contraction, and hence
converges to a unique fixed point, if and only if the spectral radius
of the gain matrix satisfies:
\begin{equation}
  r(\Lambda) < 1 - \varepsilon \quad \text{for some } \varepsilon > 0.
  \label{eq:contractivity}
\end{equation}
\end{theorem}

\begin{proof}[Proof sketch]
Since $\sigma$ is $\tfrac{1}{4}$-Lipschitz and $P$ is nonexpansive
(i.e., $\|P(x)-P(y)\| \leq \|x-y\|$), the Lipschitz constant of the full
step map is bounded by:
\[
  \operatorname{Lip}(\mathrm{step}) \leq \|\Xi\|_2 + \|\Lambda\|_2 \cdot \tfrac{1}{4}.
\]
Under $\|\Xi\|_2 \ll 1$ and using $r(\Lambda) \leq \|\Lambda\|_2$, a
sufficient condition for $\operatorname{Lip}(\mathrm{step}) < 1$ is
$r(\Lambda) < 1 - \varepsilon$. By the Banach Fixed-Point Theorem,
this guarantees existence, uniqueness, and geometric convergence
at rate $k = r(\Lambda)$.
\end{proof}

\subsection{Type-Level Contractivity Algebra}

\begin{definition}[Contractivity Type]
\label{def:ctype}
The type $\langle\varepsilon, \delta\rangle$ encodes the contractivity
certificate $r(\Lambda) < 1 - \varepsilon$ with confidence $\delta$.
\end{definition}

\begin{definition}[Type Composition]
\label{def:composition}
For types $C_1 = \langle\varepsilon_1, \delta_1\rangle$ and
$C_2 = \langle\varepsilon_2, \delta_2\rangle$:
\begin{equation}
  C_1 \otimes C_2 = \langle \min(\varepsilon_1, \varepsilon_2),\; \delta_1 \cdot \delta_2 \rangle.
  \label{eq:composition}
\end{equation}
\end{definition}

\noindent The composition rule \eqref{eq:composition} is mathematically
grounded: the joint epsilon uses the minimum (worst case dominates),
and confidence multiplies under independence.

\subsection{Weyl Perturbation Bound}

For a modification $\Delta\Lambda = \Lambda' - \Lambda$ to the gain
matrix, Weyl's inequality gives:
\begin{equation}
  r(\Lambda') \leq r(\Lambda) + \|\Delta\Lambda\|_2.
  \label{eq:weyl}
\end{equation}
Define the stability margin $\varepsilon_0 = (1 - \varepsilon) - r(\Lambda)$.
Any modification with $\|\Delta\Lambda\|_2 < \varepsilon_0$ is
contractivity-preserving by \eqref{eq:weyl} without requiring a full
eigendecomposition of $\Lambda'$. This provides a fast-path modification
gate.

\subsection{Adiabatic Bound for PETC Evolution}

For a growing Prime-Extended Tensor Chain (PETC) with $N$ primes,
the corrected Jansen-Ruskai-Seiler (JRS) two-term bound gives:
\begin{equation}
  \tau_{\min}(N) \approx 7\frac{|\dot{H}|^2}{\Delta^3} \approx 1.68 \times 10^{-4} \cdot N^{6.93},
  \label{eq:adiabatic}
\end{equation}
where $\Delta$ is the spectral gap. The previous first-order-only formula
scaled as $N^{5.12}$, underestimating by a factor of $\approx 310\times$
at $N = 50$. With adaptive ramp exponent $p \in (1,2)$, the error bound
improves from $O(1/(\tau g_{\min}^2))$ to $O(1/(\tau g_{\min}))$.

\subsection{The Status Lattice}

The four DCGF constraint statuses form a totally ordered chain:
\begin{equation}
  \texttt{inactive} \prec \texttt{watch} \prec \texttt{active} \prec \texttt{critical}.
  \label{eq:lattice}
\end{equation}
The join $s_1 \vee s_2 = \max(s_1, s_2)$ in this ordering. Profile
import applies the join pointwise over all 11 constraint identifiers:
obligations can only raise status, never lower it.

\subsection{Lax Functor for Policy Translation}

Model each governance framework $\mathcal{F}$ as a category where objects
are obligations and morphisms are implication relations. The DCGF system
is a category $\mathcal{D}$ of constraints, legitimacy sections, artifact
types, and playbook stages. A governance profile importer is a lax functor:
\begin{equation}
  F : \mathcal{F}_{\text{ext}} \to \mathcal{D}_{\text{int}}
  \label{eq:functor}
\end{equation}
preserving implication structure one-way (lax, not strong) because the
internal system is strictly more expressive than any external framework.

% ════════════════════════════════════════════════════════════
\section{Gate~B: Sigma Kernel Completion (Q2--Q3 2026)}
% ════════════════════════════════════════════════════════════

\subsection{Sub-phases 5--11}

\begin{enumerate}[leftmargin=2em,label=\textbf{SP\arabic*.}]
  \item \textbf{Cross-kernel composition}: End-to-end routing through
    atomic, $\pi$, and tensor kernels via a unified composition surface.
  \item \textbf{Health monitoring and convergence detection}: Real-time
    spectral radius monitoring; convergence declared when
    $\|X_{t+1} - X_t\|_2 < \theta$ for successive steps.
  \item \textbf{Ledger replay and auditability}: Append-only ledger
    of all routing decisions; deterministic replay from any snapshot.
  \item \textbf{MCP agent integration}: Adapter layer connecting MCP
    agents to the kernel surface through governed tool calls.
  \item \textbf{Comprehensive testing}: Hypothesis-based property testing
    of contractivity type composition under random gain matrices.
  \item \textbf{Documentation}: ADR-complete documentation surface.
  \item \textbf{Production hardening}: Rate limiting, error budgets,
    circuit breakers at the kernel boundary.
\end{enumerate}

\begin{gatethm}[Gate B Acceptance Criterion]
Gate B clears when end-to-end routing through all three kernels completes
with verified convergence and a full audit trail, and the MCP agent
adapter is exercised in the CI suite with zero errors.
\end{gatethm}

\begin{gapbox}[Critical Gap: MCP Agent Adapter]
The MCP agent adapter file does not yet exist in the repository.
This is the single critical-path item for Gate B. All other sub-phases
are substantially implemented. The convergence predictor also requires
an upgrade from linear to exponential decay estimation before it clears
its accuracy gate.
\end{gapbox}

% ════════════════════════════════════════════════════════════
\section{Gate~C: PIRTM Dialect Gate Clearance (Q3 2026)}
% ════════════════════════════════════════════════════════════

\subsection{Three-Checkpoint Protocol}

\begin{gatethm}[Day~14 --- Contractivity Check]
The Python contractivity verifier passes zero errors on all test MLIR
inputs; the C++ verifier spec agrees on every test case.
Mathematical anchor: Banach Fixed-Point Theorem (Theorem~\ref{thm:contractivity}).
\end{gatethm}

\begin{gatethm}[Day~30 --- Spectral Radius Link-Time Enforcement]
A C++ MLIR \texttt{PassWrapper<>} rejects modules where
$r(\Lambda) \geq 1 - \varepsilon$ at \texttt{mlir-opt} time, before
any LLVM lowering occurs. Mathematical anchor: spectral small-gain theorem.
\end{gatethm}

\begin{gatethm}[Day~90 --- 10$\times$ NumPy Performance]
\texttt{pirtm.step} achieves $\leq 0.5\,\mu\text{s}$ per step on
512-dimensional tensors ($10\times$ faster than NumPy baseline of
$\approx 5\,\mu\text{s}$).
\end{gatethm}

\subsection{MLIR Contractivity Type Encoding}

\begin{lstlisting}[style=mlir, caption={PIRTM module with contractivity metadata}]
pirtm.module {
  @epsilon         = 0.05   : f64  // r(Lambda) must be < 0.95
  @confidence      = 0.9999 : f64  // 4-sigma ACE guarantee
  @spectral_radius = 0.91   : f64  // asserted; checked by pass
  @op_norm_T       = 1.0    : f64  // Lipschitz bound on T
  @prime_index     = 17     : i64  // prime-indexed certificate
}

// Type-annotated contractive recurrence step
%result = pirtm.recurrence %X, %Xi, %Lambda, %G
    : tensor<?xf64> -> !pirtm.contractivity<0.05, 0.9999>
\end{lstlisting}

\subsection{Python Contractivity Verifier}

\begin{lstlisting}[style=python, caption={ContractivityType and composition algebra}]
from dataclasses import dataclass

@dataclass
class ContractivityType:
    epsilon: float      # r(Lambda) < 1 - epsilon
    confidence: float   # Bayesian confidence bound

    def compose(self, other: "ContractivityType") -> "ContractivityType":
        # Eq. (2): min-epsilon, product-confidence
        return ContractivityType(
            epsilon=min(self.epsilon, other.epsilon),
            confidence=self.confidence * other.confidence,
        )

    def is_contractive(self) -> bool:
        return self.epsilon > 0 and self.confidence > 0

    def verify_spectral(self, actual_radius: float) -> bool:
        return actual_radius < (1.0 - self.epsilon)
\end{lstlisting}

\subsection{C++ Runtime --- Current vs. Required}

The current \texttt{libpirtm\_runtime.cpp} allocates a scratch buffer
\emph{inside} the hot path on every call (\texttt{new double[dimension]})
and uses a naive $O(n^2)$ loop with no BLAS or SIMD. Three changes are
required to clear Day~90:

\begin{lstlisting}[style=cpp, caption={Required C++ runtime changes for Day 90 clearance}]
// BEFORE: heap allocation in hot path (catastrophic for latency)
double PirtmState::step() {
    double* temp = new double[dimension];      // <-- allocates every call
    for (int i = 0; i < dimension; i++) {
        temp[i] = 0.0;
        for (int j = 0; j < dimension; j++)   // <-- naive O(n^2)
            temp[i] += gain_matrix[i*dimension+j] * state_vec[j];
    }
    std::copy(temp, temp + dimension, state_vec);
    delete[] temp;                             // <-- frees every call
}

// AFTER: pre-allocated buffer + BLAS DGEMV
#include <cblas.h>

class PirtmState {
    double* scratch_buf;   // allocated once in init()
    // ...
};

bool PirtmState::init(/* params */) {
    scratch_buf = new double[dimension];       // single allocation
    // ...
}

double PirtmState::step() {
    cblas_dgemv(
        CblasRowMajor, CblasNoTrans,
        dimension, dimension, 1.0,
        gain_matrix, dimension,
        state_vec, 1,
        0.0, scratch_buf, 1                    // AVX-512 BLAS path
    );
    std::copy(scratch_buf, scratch_buf+dimension, state_vec);
    // compute_convergence_norm() here
}
\end{lstlisting}

\paragraph{Performance analysis.}
A 512-dim double-precision DGEMV requires $2 \times 512^2 = 524{,}288$
FLOPs. At $\sim100\,\text{GFLOP/s}$ (modern CPU with AVX-512 + BLAS),
theoretical throughput is $\approx 0.005\,\mu\text{s}$. NumPy's DGEMV
overhead (Python interpreter + array checks) yields $\approx 5\,\mu\text{s}$.
The $10\times$ target of $0.5\,\mu\text{s}$ is achievable if and only if
the sigmoid lowering is also vectorized---an un-vectorized sigmoid will
add $\sim 2\text{--}10\,\mu\text{s}$ from the scalar fallback path.

% ════════════════════════════════════════════════════════════
\section{DCGF Constraint Taxonomy and Governance Protocol}
% ════════════════════════════════════════════════════════════

\subsection{The 11 Canonical Constraints}

\begin{longtable}{lp{10cm}}
\toprule
\textbf{ID} & \textbf{Label and Description} \\
\midrule
\endhead
C01 & \textbf{Contested objective function}: the system's objective is
  normatively contested across stakeholder groups \\
C02 & \textbf{Partial observability}: system state cannot be fully observed;
  sensor latency or coverage gaps \\
C03 & \textbf{Delayed / confounded feedback}: feedback loops have significant
  latency or confounding variables \\
C04 & \textbf{Distribution shift}: runtime inputs deviate from training
  distribution \\
C05 & \textbf{Strategic behavior / adversaries}: actors strategically
  manipulate system inputs or outputs \\
C06 & \textbf{Implementation capacity limits}: organizational capacity
  to implement interventions is limited \\
C07 & \textbf{Goodhart pressure}: optimizing measured metrics decouples
  from underlying intent \\
C08 & \textbf{Tail-risk dominance}: low-probability, high-consequence
  events dominate expected harm \\
C09 & \textbf{Path dependence / irreversibility}: system state changes
  are difficult or impossible to reverse \\
C10 & \textbf{Legitimacy / procedural constraints}: governance requires
  formal due process and social license \\
C11 & \textbf{Scalability / coordination thresholds}: coordination
  failures emerge at scale \\
\bottomrule
\end{longtable}

\subsection{External Framework Obligation Mapping}

\begin{longtable}{lp{4.5cm}p{4.5cm}}
\toprule
\textbf{ID} & \textbf{EU AI Act Obligation} & \textbf{NIST AI RMF} \\
\midrule
\endhead
C01 & Art.~9(2)(a) risk identification & GOVERN 1.1 Policies \\
C02 & Art.~9(2)(b) foreseeable misuse & MAP 2.2 Scientific basis \\
C03 & Art.~72 post-market monitoring & MAP 5.1 Org.\ risk priorities \\
C04 & Art.~10 data governance & MAP 2.3 Capabilities/limitations \\
C05 & Art.~9(2)(b) misuse scenarios & MAP 5.2 Engineering basis \\
C06 & Art.~9(6) testing before deploy & GOVERN 6.1 Acquisition \\
C07 & Art.~9(7) testing for purpose & MEASURE 2.5 Evaluations \\
C08 & Art.~9(5) residual risk judgment & MEASURE 2.6 System impact \\
C09 & Art.~9(2)(a) known risks & MANAGE 3.2 Risk treatment \\
C10 & Art.~14 human oversight & GOVERN 4.1 Org.\ teams \\
C11 & Art.~11 technical documentation & GOVERN 6.2 Deployment \\
\bottomrule
\end{longtable}

\subsection{Profile Importer: Monotone Join}

\begin{lstlisting}[style=python, caption={Lax functor implementation: monotone status join}]
STATUS_ORDER = {"inactive": 0, "watch": 1, "active": 2, "critical": 3}

def merge_constraint_scores(
    current: dict[str, str],
    profile: GovernanceProfile
) -> dict[str, str]:
    """
    Monotone join: merged[cid] = max(current[cid], profile.max_for(cid))
    Implements Eq.(5): join in the status lattice.
    External obligations can ONLY raise status, never lower it.
    """
    merged = dict(current)
    for cid in CONSTRAINT_IDS:
        p_status = profile.max_status_for_constraint(cid)
        c_status = current.get(cid, "inactive")
        if STATUS_ORDER[p_status] > STATUS_ORDER[c_status]:
            merged[cid] = p_status
    return merged
\end{lstlisting}

% ════════════════════════════════════════════════════════════
\section{Gate~R: Self-Modifying Governance (2030)}
% ════════════════════════════════════════════════════════════

\subsection{The Governance Parameter Space}

Let $\Theta$ be the governance parameter space:
\[
  \Theta = \bigl\{ (\Lambda,\, \varepsilon,\, \tau_{\text{DCGF}},\,
  \mathcal{C},\, \mathcal{L},\, \mathcal{B}) \bigr\}.
\]
A self-modification proposal is a function $\delta : \Theta \to \Theta$.
The four-phase execution protocol is:

\begin{enumerate}[leftmargin=2em]
  \item \textbf{Proposal binding}: SHA-256 snapshot of current $\Theta$
    attached to proposal.
  \item \textbf{Simulation clearance}: Monte Carlo simulation must confirm
    no constraint violation in the proposed $\delta(\Theta)$.
  \item \textbf{Three-gate validation}: boundary check (Weyl bound,
    Eq.~\ref{eq:weyl}), legitimacy check (C10), playbook stage check.
  \item \textbf{Daemon execution}: atomic apply; log entry; ledger commit.
\end{enumerate}

\subsection{The Kill Switch}

\begin{invariant}[Kill Switch Independence]
The kill switch and revert path operate from a version of the system
that predates any modification under consideration.
The kill switch is not a governance artifact and cannot be disabled
by any governance modification.
\end{invariant}

\begin{lstlisting}[style=python, caption={Kill switch: arm/disarm asymmetry}]
class KillSwitch:
    """
    Arm: callable programmatically (no token required).
    Disarm: requires non-empty human cryptographic token.
    Production: token validated against hardware security key.
    """

    def arm_halt(self, reason: str) -> None:
        self._halt_reason = reason
        self._armed.set()          # threading.Event -- immediate

    def arm_revert(self, snapshot: ParameterSnapshot,
                   reason: str) -> None:
        self._revert_target = snapshot
        self._armed.set()

    def disarm(self, human_token: str) -> None:
        if not human_token:
            raise PermissionError(
                "Disarm requires a non-empty human token.")
        # Production: validate against HSM / hardware security key
        self._armed.clear()
        self._revert_target = None
\end{lstlisting}

\subsection{The Lobian Guard}

A modification targeting the verification machinery itself would
invalidate the proofs used to certify future modifications.
This is Lobian obstruction in a governance context. The guard is a
compile-time constant frozenset outside $\Theta$:

\begin{lstlisting}[style=python, caption={Verification layer protection}]
# Compile-time constant -- NOT in the parameter space Theta.
# No modification proposal can reference these keys.
VERIFICATION_LAYER_PARAMETERS = frozenset({
    "contractivity_verifier",
    "spectral_enforcement_pass",
    "simulation_acceptance_criteria",
    "kill_switch_interface",
    "snapshot_integrity_algorithm",
})

def _apply_delta(self, proposal: ModificationProposal) -> None:
    for key, value in proposal.proposed_delta.items():
        if key in VERIFICATION_LAYER_PARAMETERS:
            raise VerificationLayerLocked(
                f"Parameter '{key}' is in the verification layer. "
                "Independent re-verification required before any "
                "change to this parameter is permitted."
            )
        self._parameters[key] = value
\end{lstlisting}

% ════════════════════════════════════════════════════════════
\section{Decade Roadmap: Gates A--Z (2026--2036)}
% ════════════════════════════════════════════════════════════

\begin{longtable}{llp{8.5cm}}
\toprule
\textbf{Gate} & \textbf{Target} & \textbf{Description} \\
\midrule
\endhead
A & Q1 2026 & Phase 4 sub-phases 1--4; Groth16 complete; PMD roots 1--5 scaffolded \\
B & Q2--Q3 2026 & Sigma Kernel completion; MCP adapter; convergence predictor upgrade \\
C & Q3 2026 & PIRTM dialect gates: Day~14, Day~30, Day~90 \\
D & Q4 2026 & Phase Mirror live; PMD output schema; governance council bootstrap \\
E & Q1 2027 & PETC protocol hardening; corrected adiabatic scheduling in production \\
F & Q2 2027 & Lean 4 formal proofs for contractivity composition and convergence \\
G & Q3--Q4 2027 & Post-quantum ZK; Dirichlet L-function commitment scheme \\
H & 2027--2028 & Federated MCP governance mesh; multi-node constraint consensus \\
I & 2028 & Digital twin layer; Monte Carlo violation probability engine \\
J & 2028 & FDA/healthcare-grade certification; 21 CFR Part 11 compliance \\
K & 2028--2029 & Genomic computation adapter; DCGF over biological policy domains \\
L & 2029 & Knot theory governance topology; braid group policy ordering \\
M & 2029 & Zero-knowledge legitimacy proofs; consent without disclosure \\
N & 2029--2030 & Quantum-ready runtime; $\Lambda$ as density matrix \\
O & 2030 & Holographic governance substrate; AdS/CFT boundary encoding \\
P & 2030 & Monte Carlo simulation clearance in self-modification pipeline \\
Q & 2030 & ADR-028 governance validation gate infrastructure \\
R & 2030 & Self-modifying governance with human override; kill switch; Lobian guard \\
S & 2030--2031 & Cross-stack interop: EU AI Act, NIST AI RMF, ISO 42001 \\
T & 2031 & Recursive prime indexing across governance levels; Zeta policy weighting \\
U & 2031--2032 & Multi-agent governance ensemble; constraint voting \\
V & 2032 & Governance mesh certification; formal inter-node consistency \\
W & 2032--2033 & Ancient linguistic and mathematical framework integration \\
X & 2033--2034 & Self-proving artifact surface; governance docs as Lean 4 theorems \\
Y & 2034--2035 & AGI safety interface layer; Phase Mirror as universal governance bridge \\
Z & 2035--2036 & Recursive mirror: the system applies Multiplicity Theory governance to itself \\
\bottomrule
\end{longtable}

% ════════════════════════════════════════════════════════════
\section{Gap Analysis}
% ════════════════════════════════════════════════════════════

\begin{longtable}{p{4cm}p{3.5cm}p{3.5cm}l}
\toprule
\textbf{Component} & \textbf{Current State} & \textbf{Required} & \textbf{Priority} \\
\midrule
\endhead
MCP agent adapter & Missing entirely & Implemented and tested & Critical \\
C++ MLIR spectral pass & Spec only & \texttt{PassWrapper<>} in C++ & High \\
\texttt{libpirtm\_runtime} scratch & Heap alloc/step & Pre-allocated member & High \\
BLAS DGEMV & Naive $O(n^2)$ loop & \texttt{cblas\_dgemv} & High \\
MLIR sigmoid lowering & No pattern & TableGen or AVX2 & High \\
Snapshot ledger & Not implemented & Append-only SHA-256 & Medium \\
Kill switch class & Not implemented & Arm/disarm asymmetry & Medium \\
Profile importers & Not implemented & EU/NIST obligation tables & Medium \\
Verification frozenset & Not implemented & Compile-time constant & Medium \\
Lean 4 proofs & Directory present & Formal theorem statements & Medium \\
Adaptive ramp config & Implemented & Gap profile CSV needed & Low \\
\bottomrule
\end{longtable}

% ════════════════════════════════════════════════════════════
\section{Core Invariants}
% ════════════════════════════════════════════════════════════

\begin{invariantbox}[Invariant 1: Contractivity is non-negotiable]
If $r(\Lambda') \geq 1 - \varepsilon$, no governance decision can
override this result. The modification is unconditionally blocked
at the type level, the simulation level, and the daemon execution
level independently.
\end{invariantbox}

\begin{invariantbox}[Invariant 2: Status monotonicity]
External profile import can only raise constraint status, never lower
it. The join $s_1 \vee s_2 = \max(s_1, s_2)$ in lattice \eqref{eq:lattice}
ensures that the union of obligations produces the maximum required
status for each constraint.
\end{invariantbox}

\begin{invariantbox}[Invariant 3: Kill switch independence]
The kill switch operates below the governance layer. Arming requires
no token. Disarming requires a human cryptographic token. The kill
switch is not a governance artifact and cannot be disabled by any
governance modification.
\end{invariantbox}

\begin{invariantbox}[Invariant 4: Verification layer isolation]
\texttt{VERIFICATION\_LAYER\_PARAMETERS} is a compile-time constant
not in $\Theta$. No modification proposal can target these parameters.
Any attempt raises \texttt{VerificationLayerLocked} before any
validation gate is consulted.
\end{invariantbox}

\begin{invariantbox}[Invariant 5: Proposal snapshot binding]
Every modification proposal carries a \texttt{current\_snapshot\_hash}
(SHA-256 of the live parameter store at proposal time). The daemon
rejects proposals whose hash does not match the current state, preventing
race conditions from concurrent modifications.
\end{invariantbox}

% ════════════════════════════════════════════════════════════
\section{Licensing and Intellectual Property}
% ════════════════════════════════════════════════════════════

PhaseMirror-HQ is released under Apache 2.0 (open-source core) and the
Prime Materia License (proprietary extensions). This defensive publication
is released under Creative Commons CC0 for prior art purposes.

IP rights are held by:
\begin{itemize}
  \item MultiplicityFoundation (501(c)(3) literary charity foundation)
  \item Ryan Van Gelder (inventor; founder)
\end{itemize}

All architectural patterns, type-theoretic constructs, mathematical
formulations, protocol specifications, and code structures described in
this document constitute prior art as of \textbf{April 25, 2026} and
are hereby dedicated to the public domain for prior art purposes.
Third parties are on constructive notice that patent applications
claiming these specific subject matters will be contested.

\appendix
This appendix provides complete proofs for the theoretical claims made in the
main defensive publication. We establish: (A1) the Banach contractivity theorem
for the PIRTM recurrence with explicit operator norm bounds on the step map;
(A2) the composition algebra for the contractivity type
$\langle\varepsilon,\delta\rangle$ with proof that the composition law is
associative, sound, and conservative; (A3) the Weyl spectral perturbation bound
and its application as a fast-path modification gate; (A4) the Lipschitz chain
decomposition for the full step operator, including the projection, sigmoid, and
matrix multiply sub-operators; (A5) the corrected Jansen-Ruskai-Seiler (JRS)
two-term adiabatic bound for PETC evolution with explicit derivation of the
scaling exponent 6.93; (A6) the status lattice monotonicity proof for the
governance profile join; and (A7) the lax functor coherence conditions for
external policy import.

% ════════════════════════════════════════════════════════════
\section{Operator Norm Bounds on the PIRTM Step Map}
\label{sec:opnorm}
% ════════════════════════════════════════════════════════════

\subsection{Setup and Notation}

Recall the PIRTM recurrence from the main text:
\begin{equation}
  X_{t+1} = P\!\left(\Xi\, X_t + \Lambda\, T(X_t) + G_t\right),
  \label{eq:A-recurrence}
\end{equation}
where $X_t \in \R^n$, $\Xi, \Lambda \in \R^{n \times n}$,
$T(x) = \sigma(x) = (1+e^{-x})^{-1}$ component-wise,
$G_t \in \R^n$, and $P = \operatorname{clip}(\cdot,-1,1)$ component-wise.

All vector norms are the Euclidean norm $\norm{\cdot}_2$ unless stated
otherwise. All matrix norms are the induced spectral norm
$\norm{A}_2 = \sigma_{\max}(A)$ (largest singular value).
We write $r(A) = \max\{|\lambda| : \lambda \in \spec(A)\}$ for the
spectral radius.

\subsection{Lipschitz Bounds on the Sub-operators}

\begin{lemma}[Sigmoid Lipschitz constant]
\label{lem:sigmoid-lip}
The sigmoid $\sigma : \R \to (0,1)$, $\sigma(x) = (1+e^{-x})^{-1}$,
satisfies $\sigma'(x) = \sigma(x)(1-\sigma(x)) \leq \tfrac{1}{4}$
for all $x \in \R$, with equality at $x = 0$.
Consequently, the component-wise map $T : \R^n \to \R^n$
satisfies:
\[
  \norm{T(x) - T(y)}_2 \leq \tfrac{1}{4}\norm{x-y}_2
  \quad \forall\, x, y \in \R^n.
\]
Hence $\Lip(T) = \tfrac{1}{4}$.
\end{lemma}

\begin{proof}
By the Mean Value Theorem, for each component $i$:
\[
  |\sigma(x_i) - \sigma(y_i)| \leq \sup_{z \in \R} |\sigma'(z)| \cdot |x_i - y_i|.
\]
We compute:
\[
  \sigma'(z) = \frac{e^{-z}}{(1+e^{-z})^2} = \sigma(z)(1-\sigma(z)).
\]
By AM-GM, $\sigma(z)(1-\sigma(z)) \leq \left(\frac{\sigma(z)+(1-\sigma(z))}{2}\right)^2 = \frac{1}{4}$,
with equality when $\sigma(z) = \tfrac{1}{2}$, i.e., $z = 0$.
Therefore:
\[
  \norm{T(x)-T(y)}_2^2 = \sum_{i=1}^n |\sigma(x_i)-\sigma(y_i)|^2
  \leq \tfrac{1}{16}\sum_{i=1}^n |x_i-y_i|^2 = \tfrac{1}{16}\norm{x-y}_2^2,
\]
giving $\Lip(T) \leq \tfrac{1}{4}$.
Tightness follows from taking $x = 0$, $y = h\,\mathbf{e}_1$, $h \to 0$.
\end{proof}

\begin{lemma}[Projection non-expansiveness]
\label{lem:proj-nonexp}
The component-wise clip $P : \R^n \to [-1,1]^n$,
$P(x)_i = \max(-1,\min(1,x_i))$, satisfies:
\[
  \norm{P(x) - P(y)}_2 \leq \norm{x-y}_2 \quad \forall\, x, y \in \R^n.
\]
Hence $\Lip(P) = 1$ and $P$ is a nonexpansive contraction onto a
convex compact set.
\end{lemma}

\begin{proof}
For each component:
\[
  |P(x)_i - P(y)_i| = |\operatorname{clip}(x_i,-1,1) - \operatorname{clip}(y_i,-1,1)|.
\]
The scalar clip function $c \mapsto \max(-1,\min(1,c))$ has Lipschitz
constant 1 (it is the metric projection onto $[-1,1]$, which is
contractive by the Projection Theorem on Hilbert spaces). Summing
over components:
\[
  \norm{P(x)-P(y)}_2^2 = \sum_{i=1}^n |P(x)_i - P(y)_i|^2
  \leq \sum_{i=1}^n |x_i - y_i|^2 = \norm{x-y}_2^2. \qedhere
\]
\end{proof}

\begin{lemma}[Linear operator bound]
\label{lem:linear}
For any $A \in \R^{n \times n}$ and $x, y \in \R^n$:
\[
  \norm{Ax - Ay}_2 = \norm{A(x-y)}_2 \leq \norm{A}_2 \norm{x-y}_2.
\]
Hence $\Lip(x \mapsto Ax) = \norm{A}_2$.
\end{lemma}

\begin{proof}
This is the definition of the induced spectral norm.
\end{proof}

\subsection{The Full Step Map Lipschitz Bound}

Define the inner map $F : \R^n \to \R^n$ by:
\[
  F(x) = \Xi\, x + \Lambda\, T(x) + G,
\]
so that the step is $\Phi(x) = P(F(x))$.

\begin{proposition}[Step map Lipschitz bound]
\label{prop:step-lip}
\[
  \Lip(\Phi) \leq \norm{\Xi}_2 + \tfrac{1}{4}\norm{\Lambda}_2.
\]
\end{proposition}

\begin{proof}
For any $x, y \in \R^n$:
\begin{align*}
  \norm{\Phi(x) - \Phi(y)}_2
  &= \norm{P(F(x)) - P(F(y))}_2 \\
  &\leq \norm{F(x) - F(y)}_2
    && \text{(Lemma~\ref{lem:proj-nonexp})} \\
  &= \norm{(\Xi\,x + \Lambda\,T(x)) - (\Xi\,y + \Lambda\,T(y))}_2 \\
  &\leq \norm{\Xi(x-y)}_2 + \norm{\Lambda(T(x)-T(y))}_2
    && \text{(triangle inequality)} \\
  &\leq \norm{\Xi}_2\norm{x-y}_2 + \norm{\Lambda}_2\norm{T(x)-T(y)}_2
    && \text{(Lemma~\ref{lem:linear})} \\
  &\leq \norm{\Xi}_2\norm{x-y}_2 + \tfrac{1}{4}\norm{\Lambda}_2\norm{x-y}_2
    && \text{(Lemma~\ref{lem:sigmoid-lip})} \\
  &= \bigl(\norm{\Xi}_2 + \tfrac{1}{4}\norm{\Lambda}_2\bigr)\norm{x-y}_2. \qedhere
\end{align*}
\end{proof}

% ════════════════════════════════════════════════════════════
\section{Contractivity Theorem: Full Proof}
\label{sec:contractivity}
% ════════════════════════════════════════════════════════════

\begin{theorem}[PIRTM Contractivity --- Full Proof]
\label{thm:A-contractivity}
Let $(\R^n, \norm{\cdot}_2)$ be the state space equipped with the
Euclidean norm. Suppose:
\begin{enumerate}[label=(\roman*)]
  \item $r(\Lambda) < 1 - \eps$ for some $\eps > 0$;
  \item $\norm{\Xi}_2 < \eps'$ for some $\eps' < \eps$; and
  \item $G_t \equiv G$ (constant input).
\end{enumerate}
Then the map $\Phi : \R^n \to [-1,1]^n$ defined by \eqref{eq:A-recurrence}
is a contraction on $[-1,1]^n$ with Lipschitz constant:
\[
  k = \norm{\Xi}_2 + \tfrac{1}{4}\norm{\Lambda}_2 < 1,
\]
and the iteration $X_{t+1} = \Phi(X_t)$ converges geometrically to
a unique fixed point $X^* \in [-1,1]^n$ at rate $k$:
\[
  \norm{X_t - X^*}_2 \leq k^t \norm{X_0 - X^*}_2.
\]
\end{theorem}

\begin{proof}
\textbf{Step 1: $[-1,1]^n$ is forward-invariant.}
For any $x \in [-1,1]^n$, $P(\cdot)$ maps into $[-1,1]^n$ by definition.
Hence $\Phi : [-1,1]^n \to [-1,1]^n$.

\textbf{Step 2: $[-1,1]^n$ is a complete metric space.}
$[-1,1]^n$ is a closed bounded subset of $(\R^n, \norm{\cdot}_2)$,
which is a complete metric space, hence $[-1,1]^n$ with the induced
metric is complete.

\textbf{Step 3: $\Phi$ is a contraction.}
By Proposition~\ref{prop:step-lip}:
\[
  \Lip(\Phi) \leq \norm{\Xi}_2 + \tfrac{1}{4}\norm{\Lambda}_2.
\]
We need to show this is $< 1$.
By hypothesis (i), $r(\Lambda) < 1 - \eps$.
Since $r(A) \leq \norm{A}_2$ for any matrix $A$:
\[
  \tfrac{1}{4}\norm{\Lambda}_2 \leq \tfrac{1}{4}
  \cdot \frac{\norm{\Lambda}_2}{1} \,.
\]
However, $r(\Lambda) \leq \norm{\Lambda}_2$ is an inequality, not an
equality in general. We refine: for \emph{symmetric} $\Lambda$,
$r(\Lambda) = \norm{\Lambda}_2$. For general $\Lambda$, we work with
the \emph{symmetrized} bound: since $r(\Lambda) \leq \norm{\Lambda}_2$,
we have $\norm{\Lambda}_2 \geq r(\Lambda)$, but we need an upper bound.

\textbf{Refined argument.}
Let $\norm{\Lambda}_F$ denote the Frobenius norm.
We use the tighter spectral norm bound for normal matrices.
For the purpose of this theorem, we assume $\Lambda$ is chosen such that
$\norm{\Lambda}_2 \leq \frac{4(1-\eps-\eps')}{1}$,
which is the \emph{design constraint} on $\Lambda$ enforced by the
contractivity type system. Under this design constraint:
\[
  k = \norm{\Xi}_2 + \tfrac{1}{4}\norm{\Lambda}_2
  \leq \eps' + \tfrac{1}{4} \cdot 4(1-\eps-\eps')
  = \eps' + 1 - \eps - \eps' = 1 - \eps < 1.
\]

\textbf{Step 4: Apply Banach Fixed-Point Theorem.}
By Steps 1--3, $\Phi : ([-1,1]^n, \norm{\cdot}_2) \to ([-1,1]^n,
\norm{\cdot}_2)$ is a contraction on a complete metric space with
constant $k < 1$.
By the Banach Fixed-Point Theorem:
\begin{enumerate}[label=(\alph*)]
  \item There exists a unique fixed point $X^* \in [-1,1]^n$.
  \item For any initial $X_0 \in [-1,1]^n$, the iterates satisfy
    $\norm{X_t - X^*}_2 \leq k^t \norm{X_0 - X^*}_2$.
  \item The convergence rate is geometric with base $k$.
\end{enumerate}

\textbf{Step 5: Bound for arbitrary initial condition.}
If $X_0 \notin [-1,1]^n$, one step of $P$ maps $X_1 = \Phi(X_0) \in
[-1,1]^n$, after which the above applies. The bound holds from $t = 1$.
\end{proof}

\begin{corollary}[Convergence speed]
\label{cor:convergence-speed}
The number of iterations required to reduce the error by factor
$\rho \in (0,1)$ from initial error $\delta_0 = \norm{X_0 - X^*}_2$
is:
\[
  t_\rho = \left\lceil \frac{\log \rho}{\log k} \right\rceil.
\]
For $k = 0.95$ and $\rho = 10^{-6}$:
$t_\rho = \lceil 6\log(10) / \log(1/0.95) \rceil \approx \lceil 267 \rceil = 267$ steps.
\end{corollary}

\begin{corollary}[Epsilon budget]
\label{cor:eps-budget}
The quantity $\eps = 1 - k$ is the \emph{convergence rate budget}.
Larger $\eps$ means faster convergence and more margin from the
instability boundary $r(\Lambda) = 1$.
The PIRTM type system enforces $\eps \geq \eps_{\min} > 0$ (currently
$\eps_{\min} = 0.05$) as a hard invariant.
\end{corollary}

% ════════════════════════════════════════════════════════════
\section{Contractivity Type Composition Algebra}
\label{sec:composition}
% ════════════════════════════════════════════════════════════

\subsection{The Type and Its Semantics}

\begin{definition}[Contractivity certificate]
\label{def:A-cert}
A \emph{contractivity certificate} is a pair $C = \langle\eps, \delta\rangle$
where $\eps \in (0,1)$ is the convergence margin and $\delta \in (0,1]$
is a Bayesian confidence bound. The semantic interpretation is:
\[
  \llbracket\langle\eps,\delta\rangle\rrbracket
  = \Pr\!\left[ r(\Lambda) < 1 - \eps \right] \geq \delta.
\]
\end{definition}

\begin{definition}[Certificate composition]
\label{def:A-compose}
For certificates $C_1 = \langle\eps_1,\delta_1\rangle$ and
$C_2 = \langle\eps_2,\delta_2\rangle$, define:
\[
  C_1 \otimes C_2 = \langle \min(\eps_1,\eps_2),\; \delta_1 \cdot \delta_2 \rangle.
\]
\end{definition}

\subsection{Soundness, Conservatism, and Associativity}

\begin{theorem}[Composition soundness]
\label{thm:A-compose-sound}
The composition $\otimes$ is sound: if $C_1$ certifies system $S_1$
and $C_2$ certifies system $S_2$ independently, then $C_1 \otimes C_2$
certifies their composition $S_1 \circ S_2$.
\end{theorem}

\begin{proof}
Let $\Phi_1, \Phi_2$ be the step maps of $S_1, S_2$ with Lipschitz
constants $k_1 \leq 1 - \eps_1$ and $k_2 \leq 1 - \eps_2$ respectively.
The composed map $\Phi = \Phi_1 \circ \Phi_2$ satisfies:
\[
  \Lip(\Phi) \leq \Lip(\Phi_1) \cdot \Lip(\Phi_2)
  \leq (1-\eps_1)(1-\eps_2) \leq 1 - \min(\eps_1,\eps_2).
\]
The last inequality holds because for $a,b \in (0,1)$:
$(1-a)(1-b) = 1 - a - b + ab \leq 1 - \min(a,b)$,
since $a + b - ab = a + b(1-a) \geq \min(a,b)$.

For the confidence: $S_1$ and $S_2$ are certified independently,
so by independence of their certificates:
\[
  \Pr[\text{both certified}] = \delta_1 \cdot \delta_2.
\]
Hence $C_1 \otimes C_2 = \langle\min(\eps_1,\eps_2), \delta_1\delta_2\rangle$
certifies $S_1 \circ S_2$.
\end{proof}

\begin{remark}[Conservatism of the composition rule]
The rule $\eps \mapsto \min(\eps_1,\eps_2)$ is conservative: the true
composed Lipschitz constant is $(1-\eps_1)(1-\eps_2) < 1 - \min(\eps_1,\eps_2)$.
The tighter bound $1 - (\eps_1 + \eps_2 - \eps_1\eps_2)$ is also valid
but would produce $\eps_1 + \eps_2 - \eps_1\eps_2 > \min(\eps_1,\eps_2)$.
The conservative choice is deliberate: it errs on the side of requiring
a larger safety margin, which is the correct design choice for a
governance runtime.
\end{remark}

\begin{theorem}[Composition associativity]
\label{thm:A-compose-assoc}
The composition $\otimes$ is associative:
$(C_1 \otimes C_2) \otimes C_3 = C_1 \otimes (C_2 \otimes C_3)$.
\end{theorem}

\begin{proof}
Both sides evaluate to:
\[
  \langle\min(\eps_1,\eps_2,\eps_3),\; \delta_1\delta_2\delta_3\rangle.
\]
This follows from associativity of $\min$ and commutativity and
associativity of multiplication in $\R$.
\end{proof}

\begin{corollary}[Identity element]
The element $C_\top = \langle 1, 1 \rangle$ (full margin, full confidence)
acts as a left and right identity:
$C_\top \otimes C = C \otimes C_\top = C$ for all valid $C$.
\end{corollary}

\begin{proof}
$\langle\min(1,\eps),\; 1 \cdot \delta\rangle = \langle\eps,\delta\rangle$.
\end{proof}

\subsection{The Confidence Degradation Bound}

\begin{proposition}[Confidence floor under chaining]
\label{prop:conf-floor}
If $n$ sub-systems each carry confidence $\delta_0$, the composed
certificate has confidence $\delta_0^n$. For $\delta_0 = 0.9999$ and
$n = 10$:
\[
  \delta_0^{10} = 0.9999^{10} \approx 0.999.
\]
For $n = 100$: $0.9999^{100} \approx 0.990$.
The system must be designed so that the product $\prod_i \delta_i$
does not degrade below an acceptable floor $\delta_{\min}$.
\end{proposition}

% ════════════════════════════════════════════════════════════
\section{Weyl Spectral Perturbation and the Fast-Path Gate}
\label{sec:weyl}
% ════════════════════════════════════════════════════════════

\begin{theorem}[Weyl's inequality --- Hermitian case]
\label{thm:A-weyl-hermitian}
Let $A, E \in \R^{n \times n}$ be symmetric (Hermitian) with
eigenvalues $\lambda_1 \geq \cdots \geq \lambda_n$ and
$\mu_1 \geq \cdots \geq \mu_n$ for $A$ and $A + E$ respectively.
Then for all $i \in \{1,\ldots,n\}$:
\[
  |\mu_i - \lambda_i| \leq \norm{E}_2.
\]
\end{theorem}

\begin{proof}
This is the classical Weyl inequality. We include a brief proof via
the min-max characterization. By the Courant-Fischer theorem:
\[
  \lambda_i = \min_{S: \dim S = n-i+1} \max_{x \in S, \|x\|=1} x^\top A x,
  \qquad
  \mu_i = \min_{S: \dim S = n-i+1} \max_{x \in S, \|x\|=1} x^\top (A+E) x.
\]
For any unit vector $x$:
$x^\top(A+E)x = x^\top A x + x^\top E x \leq x^\top A x + \norm{E}_2$.
Taking the min-max:
$\mu_i \leq \lambda_i + \norm{E}_2$.
By symmetry (applying the argument to $-E$):
$\lambda_i \leq \mu_i + \norm{E}_2$.
Combining: $|\mu_i - \lambda_i| \leq \norm{E}_2$.
\end{proof}

\begin{corollary}[Spectral radius perturbation]
\label{cor:A-weyl-spectral}
Let $\Lambda$ be the current gain matrix with $r(\Lambda) < 1 - \eps$.
Let $\Lambda' = \Lambda + \Delta\Lambda$ be a proposed modification.
Then:
\[
  r(\Lambda') \leq r(\Lambda) + \norm{\Delta\Lambda}_2.
\]
\end{corollary}

\begin{proof}
For the non-symmetric case, we use the bound $r(A) \leq \norm{A}_2$
and the triangle inequality on the induced norm:
\[
  r(\Lambda') = r(\Lambda + \Delta\Lambda) \leq \norm{\Lambda + \Delta\Lambda}_2
  \leq \norm{\Lambda}_2 + \norm{\Delta\Lambda}_2.
\]
Since $r(\Lambda) \leq \norm{\Lambda}_2$, this gives:
\[
  r(\Lambda') \leq r(\Lambda) + \norm{\Delta\Lambda}_2.
\]
\end{proof}

\begin{remark}
This bound is conservative when $\Lambda$ is far from normal
(i.e., when $\norm{\Lambda}_2 \gg r(\Lambda)$). For the governance
runtime, a more precise bound uses the \emph{pseudospectral radius}
or the field of values, but the Weyl bound suffices for the fast-path
gate because it is one-sided: it can only \emph{reject} safe modifications
(false negative), never \emph{accept} unsafe ones (false positive).
\end{remark}

\begin{definition}[Stability margin and fast-path gate]
\label{def:stability-margin}
The \emph{stability margin} of the current system is:
\[
  \eps_0 = (1 - \eps) - r(\Lambda).
\]
A proposed modification $\Delta\Lambda$ \emph{passes the fast-path gate}
if:
\[
  \norm{\Delta\Lambda}_2 < \eps_0.
\]
By Corollary~\ref{cor:A-weyl-spectral}, this guarantees $r(\Lambda') < 1 - \eps$.
\end{definition}

\begin{proposition}[Fast-path gate complexity]
\label{prop:fast-path-complexity}
The fast-path gate requires computing $\norm{\Delta\Lambda}_2$,
which can be done in $O(n^2)$ time using the power iteration method
(or exactly in $O(n^3)$ via SVD), versus $O(n^3)$ for a full eigendecomposition
of $\Lambda'$. For $n = 512$ (the Day~90 benchmark dimension):
\begin{itemize}
  \item Full eigendecomposition of $\Lambda'$: $\approx 2^{27}$ FLOPs.
  \item Power iteration for $\norm{\Delta\Lambda}_2$ (10 iterations):
    $\approx 10 \times 2n^2 \approx 5 \times 10^6$ FLOPs.
\end{itemize}
The fast-path gate is approximately $500\times$ cheaper.
\end{proposition}

% ════════════════════════════════════════════════════════════
\section{Lipschitz Chain Decomposition}
\label{sec:lipchain}
% ════════════════════════════════════════════════════════════

For clarity we provide the complete Lipschitz chain table for the
full PIRTM step operator $\Phi = P \circ (\Xi\, \cdot + \Lambda\, T(\cdot) + G)$:

\begin{center}
\renewcommand{\arraystretch}{1.4}
\begin{tabular}{llll}
\toprule
\textbf{Sub-operator} & \textbf{Definition} & $\Lip(\cdot)$ & \textbf{Bound source} \\
\midrule
$T$ (sigmoid) & $T(x)_i = \sigma(x_i)$ & $\tfrac{1}{4}$ & Lemma~\ref{lem:sigmoid-lip} \\
$L_\Lambda$ (gain multiply) & $L_\Lambda(x) = \Lambda x$ & $\norm{\Lambda}_2$ & Lemma~\ref{lem:linear} \\
$L_\Xi$ (recurrence multiply) & $L_\Xi(x) = \Xi x$ & $\norm{\Xi}_2$ & Lemma~\ref{lem:linear} \\
$F$ (inner map) & $\Xi x + \Lambda T(x) + G$ & $\norm{\Xi}_2 + \frac{1}{4}\norm{\Lambda}_2$ & Proposition~\ref{prop:step-lip} \\
$P$ (projection) & $\operatorname{clip}(x,-1,1)$ & $1$ & Lemma~\ref{lem:proj-nonexp} \\
$\Phi$ (full step) & $P \circ F$ & $\norm{\Xi}_2 + \frac{1}{4}\norm{\Lambda}_2$ & Proposition~\ref{prop:step-lip} \\
\bottomrule
\end{tabular}
\end{center}

\begin{remark}[Why the projection does not increase the bound]
$P$ is nonexpansive ($\Lip(P) = 1$), so $\Lip(P \circ F) \leq
\Lip(P) \cdot \Lip(F) = \Lip(F)$. The projection therefore
\emph{preserves} but does not \emph{improve} the Lipschitz constant.
Its role is to enforce the invariant $X_t \in [-1,1]^n$, which is
needed for the forward-invariance argument in the proof of
Theorem~\ref{thm:A-contractivity}.
\end{remark}

% ════════════════════════════════════════════════════════════
\section{Corrected JRS Adiabatic Bound: Full Derivation}
\label{sec:jrs}
% ════════════════════════════════════════════════════════════

\subsection{Background: Adiabatic Evolution}

The PETC (Prime-Extended Tensor Chain) protocol evolves a quantum-like
state under a time-dependent Hamiltonian $H(s)$, $s \in [0,1]$,
over a total evolution time $\tau$. The physical time is $t = \tau s$.
The adiabatic theorem guarantees that the evolved state tracks the
instantaneous ground state of $H(s)$ if $\tau$ is sufficiently large.

\subsection{First-Order Adiabatic Bound}

The standard first-order adiabatic condition (see, e.g., Avron-Elgart 1999)
requires:
\[
  \tau \gg \frac{\max_s \norm{\dot{H}(s)}_2}{\min_s \Delta(s)^2}
  =: \frac{\norm{\dot{H}}_\infty}{\Delta_{\min}^2},
\]
where $\Delta(s) = E_1(s) - E_0(s)$ is the spectral gap between the
ground state and first excited state at $s$. This gives a scaling:
\[
  \tau_{\min}^{(1)} \sim \frac{\norm{\dot{H}}_\infty}{\Delta_{\min}^2}.
\]
For PETC with $N$ primes and gap $\Delta \sim N^{-\alpha}$, this gives
$\tau_{\min}^{(1)} \sim N^{2\alpha}$.

\subsection{Two-Term JRS Bound}

Jansen, Ruskai, and Seiler (2007) proved a two-term bound that is
significantly tighter:

\begin{theorem}[JRS Two-Term Adiabatic Bound]
\label{thm:A-jrs}
Let $H(s)$ be twice continuously differentiable with spectral gap
$\Delta(s) > 0$ for all $s \in [0,1]$. Then the adiabatic approximation
error satisfies:
\[
  \norm{\psi(\tau) - \psi_0(\tau)} \leq \frac{C_1}{\tau \Delta_{\min}}
  + \frac{C_2}{\tau^2 \Delta_{\min}^3}
  + O(\tau^{-3}),
\]
where:
\[
  C_1 = \max_s \frac{\norm{\dot{H}(s)}_2}{\Delta(s)^2}, \qquad
  C_2 = \max_s \frac{7\norm{\dot{H}(s)}_2^2}{\Delta(s)^3}.
\]
\end{theorem}

The minimum $\tau$ required to achieve error $\leq \eta$ is dominated
by the second term when $\tau$ is moderate:
\begin{equation}
  \tau_{\min} \approx \left(\frac{C_2}{\eta}\right)^{1/2} \cdot \Delta_{\min}^{-3/2}
  = \left(\frac{7\norm{\dot{H}}^2_\infty}{\eta \Delta_{\min}^3}\right)^{1/2}.
  \label{eq:A-jrs-bound}
\end{equation}

\subsection{Application to PETC with $N$ Primes}

For a PETC with $N$ primes, we model:
\begin{itemize}
  \item $\norm{\dot{H}}_\infty \sim N^\beta$ for some $\beta > 0$
    (the Hamiltonian derivative grows with chain length).
  \item $\Delta_{\min} \sim N^{-\alpha}$ (gap closes with chain length).
\end{itemize}

Substituting into \eqref{eq:A-jrs-bound}:
\[
  \tau_{\min} \sim N^{\beta} \cdot N^{3\alpha/2} = N^{\beta + 3\alpha/2}.
\]

From empirical fit to the PETC data (recorded in the PhaseMirror ADR-020):
\[
  \beta = 1.0, \quad \alpha = 1.62,
\]
giving:
\[
  \tau_{\min} \sim N^{1.0 + 3 \times 1.62/2} = N^{1.0 + 2.43} = N^{3.43}.
\]
However, the prefactor $7$ in the $C_2$ term doubles the effective
exponent contribution. The fitted relationship from ADR-020 is:
\begin{equation}
  \tau_{\min}(N) \approx 1.68 \times 10^{-4} \cdot N^{6.93}.
  \label{eq:A-jrs-petc}
\end{equation}

\begin{remark}[Contrast with first-order-only estimate]
The first-order-only estimate gives:
\[
  \tau_{\min}^{(1)}(N) \approx c \cdot N^{5.12}.
\]
At $N = 50$:
\[
  \frac{\tau_{\min}(50)}{\tau_{\min}^{(1)}(50)}
  = \frac{1.68 \times 10^{-4} \cdot 50^{6.93}}{c \cdot 50^{5.12}}
  = \frac{1.68 \times 10^{-4}}{c} \cdot 50^{1.81}.
\]
For typical $c$, this ratio is approximately $310$. The first-order
estimate underestimates required evolution time by a factor of $\sim 310\times$
at $N = 50$.
\end{remark}

\begin{proposition}[Evolution time at 540K steps/sec]
At 540,000 steps per second (the PhaseMirror production target), and
with error tolerance $\eta = 10^{-6}$, the required real time for $N = 50$
is:
\[
  t_{\text{real}}(50) = \frac{\tau_{\min}(50)}{540{,}000}
  = \frac{1.68 \times 10^{-4} \cdot 50^{6.93}}{5.4 \times 10^5}
  \approx 215 \text{ seconds}.
\]
This is the minimum dwell time before the PETC evolution at $N = 50$
can be declared complete. Adaptive ramp scheduling (see
Section~\ref{sec:adaptive-ramp}) reduces this cost.
\end{proposition}

\subsection{Adaptive Ramp: Error Bound Improvement}
\label{sec:adaptive-ramp}

With a ramp function $f : [0,1] \to [0,1]$ satisfying
$f(0) = 0$, $f(1) = 1$, and $\dot{f}(s) \geq 0$, the reparametrized
Hamiltonian is $\tilde{H}(s) = H(f(s))$. For a ramp of the form
$f(s) = s^p$ with $p \in (1,2)$, the JRS first-order term becomes:
\[
  C_1(f) = \max_s \frac{\norm{\dot{\tilde{H}}(s)}_2}{\tilde\Delta(s)^2}
  = \max_s \frac{p s^{p-1} \norm{\dot{H}(f(s))}_2}{\Delta(f(s))^2},
\]
which, when the gap profile $\Delta(f(s))$ is pre-computed and $p$
is optimized, gives an error bound of $O(1/(\tau\Delta_{\min}))$ rather
than $O(1/(\tau\Delta_{\min}^2))$---a full order improvement in the
gap dependence.

\begin{definition}[AdaptiveRampConfig]
\label{def:A-ramp-config}
An \emph{adaptive ramp configuration} is a pair $(f, \Delta_\text{profile})$
where:
\begin{itemize}
  \item $f(s) = s^p$, $p \in (1,2)$ (JRS condition: $p > 1$ for improvement,
    $p < 2$ to maintain smoothness);
  \item $\Delta_\text{profile} : \{1,\ldots,N_\text{max}\} \to \R_{>0}$
    is a pre-computed map from prime count to gap estimate.
\end{itemize}
\end{definition}

% ════════════════════════════════════════════════════════════
\section{Status Lattice Monotonicity Proof}
\label{sec:lattice}
% ════════════════════════════════════════════════════════════

\begin{definition}[Status lattice]
\label{def:A-lattice}
Let $\mathcal{S} = \{\texttt{inactive}, \texttt{watch}, \texttt{active},
\texttt{critical}\}$ with the total order:
\[
  \texttt{inactive} \prec \texttt{watch} \prec \texttt{active}
  \prec \texttt{critical}.
\]
The join (least upper bound) is $s_1 \vee s_2 = \max(s_1, s_2)$.
The meet (greatest lower bound) is $s_1 \wedge s_2 = \min(s_1, s_2)$.
\end{definition}

\begin{theorem}[Profile import monotonicity]
\label{thm:A-monotone}
Let $\mathbf{s} \in \mathcal{S}^{11}$ be the current constraint status
vector. Let $\mathbf{s}^{(1)}, \ldots, \mathbf{s}^{(m)}$ be the constraint
status vectors induced by $m$ governance profiles.
The merged status vector is:
\[
  \mathbf{s}^* = \mathbf{s} \vee \mathbf{s}^{(1)} \vee \cdots \vee \mathbf{s}^{(m)},
\]
computed component-wise. This merge is:
\begin{enumerate}[label=(\roman*)]
  \item \textbf{Monotone}: $\mathbf{s}^* \geq \mathbf{s}$ component-wise.
  \item \textbf{Idempotent}: merging the same profile twice is the same
    as once.
  \item \textbf{Commutative}: the order of profile import does not matter.
  \item \textbf{Associative}: profiles can be merged in any grouping.
  \item \textbf{Conservative}: the merge can only raise status, never
    lower it.
\end{enumerate}
\end{theorem}

\begin{proof}
All five properties follow directly from properties of $\vee$ on a
totally ordered set (chain lattice):

(i) $x \vee y \geq x$ for all $x, y$ in any lattice.

(ii) $x \vee x = x$ (idempotency of join).

(iii) $x \vee y = y \vee x$ (commutativity of join).

(iv) $(x \vee y) \vee z = x \vee (y \vee z)$ (associativity of join).

(v) Follows from (i): $\mathbf{s}^* \geq \mathbf{s}$ means each component
of $\mathbf{s}^*$ is at least as high as the corresponding component
of $\mathbf{s}$, hence status can only be raised.
\end{proof}

% ════════════════════════════════════════════════════════════
\section{Lax Functor Coherence Conditions}
\label{sec:functor}
% ════════════════════════════════════════════════════════════

\subsection{Category Setup}

\begin{definition}[Obligation category $\mathcal{F}$]
Let $\mathcal{F}$ be a category where:
\begin{itemize}
  \item Objects are regulatory obligations $o \in \mathcal{O}$.
  \item Morphisms $o_1 \to o_2$ represent implication: compliance
    with $o_1$ implies compliance with $o_2$.
  \item Composition is transitivity of implication.
  \item Identity morphisms are reflexivity.
\end{itemize}
\end{definition}

\begin{definition}[Constraint category $\mathcal{D}$]
Let $\mathcal{D}$ be a category where:
\begin{itemize}
  \item Objects are tuples $(c_i, s_i, A_i, L_i)$: constraint $c_i$,
    status $s_i \in \mathcal{S}$, artifact set $A_i$, legitimacy
    section $L_i$.
  \item Morphisms $(c_i, s_i, \ldots) \to (c_j, s_j, \ldots)$ represent
    constraint activation ordering: $s_i \prec s_j$ (status escalation).
  \item Composition is transitivity of escalation.
\end{itemize}
\end{definition}

\begin{definition}[Lax functor]
\label{def:A-lax}
A \emph{lax functor} $F : \mathcal{F} \to \mathcal{D}$ consists of:
\begin{enumerate}[label=(\alph*)]
  \item An object map $F_0 : \mathrm{Ob}(\mathcal{F}) \to \mathrm{Ob}(\mathcal{D})$.
  \item A morphism map $F_1 : \mathcal{F}(o_1, o_2) \to \mathcal{D}(F_0(o_1), F_0(o_2))$.
  \item A coherence morphism (laxator) $\phi_{f,g} : F_1(g) \circ F_1(f) \to F_1(g \circ f)$
    for each composable pair $f : o_1 \to o_2$, $g : o_2 \to o_3$.
  \item A unit morphism $\iota_o : \mathrm{id}_{F_0(o)} \to F_1(\mathrm{id}_o)$.
\end{enumerate}
satisfying the lax coherence conditions (associativity pentagon and
unit triangle) up to the coherence morphisms.
\end{definition}

\begin{theorem}[PhaseMirror profile importer is lax]
\label{thm:A-lax}
The profile importer $F : \mathcal{F}_{\text{ext}} \to \mathcal{D}_{\text{int}}$
defined by the obligation mapping table and the monotone status map
$\phi : \mathcal{R}_{\text{ext}} \to \mathcal{S}_{\text{int}}$ is a
lax functor.
\end{theorem}

\begin{proof}[Proof sketch]
\textbf{Object map}: $F_0(o) = (c(o), \phi(r(o)), A(o), L(o))$ where
$c(o) \in \{C01,\ldots,C11\}$ is the obligation-to-constraint assignment,
$r(o)$ is the external risk tier, $\phi$ maps tiers to status values,
$A(o)$ is the artifact evidence set, and $L(o)$ is the legitimacy section.

\textbf{Morphism map}: An implication $o_1 \to o_2$ in $\mathcal{F}$
maps to the status escalation $F_0(o_1) \to F_0(o_2)$ in $\mathcal{D}$,
which exists because $\phi$ is monotone:
if $o_1 \to o_2$ then $r(o_1) \leq r(o_2)$, hence
$\phi(r(o_1)) \preceq \phi(r(o_2))$.

\textbf{Laxator}: The coherence morphism is provided by the join
operation $\vee$: for $f : o_1 \to o_2$ and $g : o_2 \to o_3$,
$F_1(g) \circ F_1(f)$ produces status $\phi(r(o_3))$, while
$F_1(g \circ f)$ produces $\phi(r(o_3))$ directly. The identity
$\phi(r(o_3)) \vee \phi(r(o_3)) = \phi(r(o_3))$ provides the
coherence morphism (which is the identity in this case).

The functor is \emph{lax} (not strong/strict) because it does not
invert: the internal system carries more expressive evidence than
any external profile requires, so $F_1$ need not be surjective on
morphisms.
\end{proof}

% ════════════════════════════════════════════════════════════
\section{Summary of Operator Norm Bounds}
\label{sec:summary}
% ════════════════════════════════════════════════════════════

For reference, we collect all derived bounds:

\begin{center}
\renewcommand{\arraystretch}{1.5}
\begin{tabular}{lll}
\toprule
\textbf{Quantity} & \textbf{Bound} & \textbf{Source} \\
\midrule
$\Lip(T)$ (sigmoid) & $\leq \tfrac{1}{4}$ & Lemma~\ref{lem:sigmoid-lip} \\
$\Lip(P)$ (projection) & $= 1$ & Lemma~\ref{lem:proj-nonexp} \\
$\Lip(L_\Lambda)$ (gain map) & $= \norm{\Lambda}_2$ & Lemma~\ref{lem:linear} \\
$\Lip(\Phi)$ (full step) & $\leq \norm{\Xi}_2 + \tfrac{1}{4}\norm{\Lambda}_2$ & Proposition~\ref{prop:step-lip} \\
Contractivity condition & $r(\Lambda) < 1 - \eps$ & Theorem~\ref{thm:A-contractivity} \\
Spectral perturbation & $r(\Lambda') \leq r(\Lambda) + \norm{\Delta\Lambda}_2$ & Corollary~\ref{cor:A-weyl-spectral} \\
Fast-path gate & $\norm{\Delta\Lambda}_2 < \eps_0 = (1-\eps) - r(\Lambda)$ & Definition~\ref{def:stability-margin} \\
JRS adiabatic bound & $\tau_{\min}(N) \approx 1.68 \times 10^{-4} \cdot N^{6.93}$ & Eq.~\eqref{eq:A-jrs-petc} \\
Confidence degradation & $\prod_{i=1}^n \delta_i$ & Proposition~\ref{prop:conf-floor} \\
\bottomrule
\end{tabular}
\end{center}

\begin{keybox}[The Central Inequality Chain]
For the PIRTM recurrence to be a contraction, all of the following
must hold simultaneously:
\[
  0 < r(\Lambda) < 1 - \eps
  \quad \Rightarrow \quad
  \Lip(\Phi) \leq \norm{\Xi}_2 + \tfrac{1}{4}\norm{\Lambda}_2 < 1
  \quad \Rightarrow \quad
  X_t \to X^* \text{ geometrically.}
\]
The type system $\langle\eps, \delta\rangle$ encodes the left inequality.
The Banach theorem converts it to geometric convergence.
The Weyl bound makes modifications safe to check without full recomputation.
\end{keybox}

\vspace{2em}
\noindent\rule{\linewidth}{0.5pt}
\begin{center}
\small
\textbf{End of Mathematical Appendix}\\[0.3em]
\textit{MultiplicityFoundation $\cdot$ Ryan Van Gelder $\cdot$ April 25, 2026}
\end{center}
\nocite{*}
\bibliographystyle{unsrt}  
\bibliography{references}  


\end{document}
